\lecture{11}{29. September 2025}{Stress Design and Deformation}

\subsection{Average Shear stress}
Similarly it can be shown that the average shear stress $\tau_{\mathrm{avg}}$ is given as
\[ 
\tau_{\mathrm{avg}} = \frac{V}{A}
\]
where, $V$ is the resultant internal shear force determined from the equations of equilibrium and $A$ is the area of the section. 


\subsection{Allowable Stress Design}
To ensure safety it is needed to restrict the allowed load on a structure to one that is less than what the structure can theoretically support. One method for doing this is to use the \textit{factor of safety} ($\mathrm{F.S.}$) which is a ratio between the failure load $F_{\mathrm{fail}}$ to the allowable load $F_{\mathrm{allow}}$,
\[ 
\mathrm{F.S.} = \frac{F_{\text{fail}}}{F_{\text{allow}}}
.\]
Here $F_{\text{fail}}$ is typically found through experimental testing of the material. 


\subsection{Strain}
The normal strain has been defined in materials science as,
\[ 
\epsilon_{\mathrm{avg}} = \frac{\Delta L}{L}
.\]
We recall that this is a dimensionless quantity.

Deformation, however, not only presents as normal strain. I.e. line segments of elements will not only expand and contract but also change direction. The \textit{change in angle} between two segments, originally perpendicular, is referred to as the \textit{shear strain}, denoted by $\gamma$ and always measured in \unit{rad}. 
